\subsection{Auftragsbasierte Prozesse}
\subsubsection{Grundprinzip der auftragsbasierten Lagerung}
Die auftragsbasierte Lagerung ermöglicht es dem Benutzer, den Lagerprozess nicht mehr durch direkte Eingaben, sondern über definierte Aufträge zu steuern.\\
Im Falle der manuellen Lagerung gibt der Benutzer seinen gewünschten Lagerplatz vor. Alternativ besteht die Möglichkeit, ein bestimmtes Bauteil auszuwählen, woraufhin der entsprechende Prozess initiiert wird.\\
Bei der auftragsbasierten Lagerung wird der Lagerprozess hingegen über ein \textbf{Eingabegerät} gesteuert. Dabei erfolgt die Eingabe nicht mehr positionsorientiert, sondern auftragsorientiert.\\

Ein Auftrag enthält die notwendigen Informationen zur Identifikation eines Lagergutes, beispielsweise eine Artikel- oder Bauteilnummer. Auf Grundlage dieser Informationen übernimmt das System die weitere Verarbeitung und trifft eigenständig Entscheidungen hinsichtlich der Lagerbewegung.\\
In vielen Fällen werden diese Informationen über eine Datenbank bereitgestellt, wodurch eine flexible und modulare Anpassung der Auftragsstruktur jederzeit möglich ist.\\

Die logische Verarbeitung der Aufträge erfolgt dabei softwareseitig, während die physische Umsetzung weiterhin durch die speicherprogrammierbare Steuerung (SPS) realisiert wird. Dadurch entsteht eine klare Trennung zwischen Entscheidungslogik und Anlagensteuerung.\\

Der Ablauf eines solchen Prozesses beginnt mit der Erfassung eines Auftrags, beispielsweise durch das Einlesen eines Codes mittels Scanner. Anschließend wird dieser Auftrag verarbeitet und die entsprechenden Lageroperationen werden automatisch ausgelöst.\\

Ziel dieser Vorgehensweise ist es, manuelle Eingriffe zu reduzieren, die Prozesseffizienz zu steigern sowie eine höhere Nachvollziehbarkeit und Fehlerreduktion innerhalb der Lagerlogistik zu gewährleisten.
\subsubsection{Ablauf der auftragsbasierten Lagerung}

Der Ablauf der auftragsbasierten Lagerung basiert auf dem Zusammenspiel mehrerer Systemkomponenten, welche unterschiedlichen Ebenen zugeordnet sind\cite{Grundlagen_Logistik}. Diese Ebenen umfassen sowohl die Eingabe- und Verarbeitungsebene als auch die Steuerungs- und Ausführungsebene der Anlage. Ziel dieser Struktur ist es, die einzelnen Prozessschritte klar voneinander zu trennen und die Datenflüsse innerhalb des Systems transparent darzustellen.\\

Im Gegensatz zu klassischen Lagerprozessen erfolgt die Steuerung hierbei nicht direkt durch den Bediener, sondern durch eine softwaregestützte Verarbeitung von Auftragsdaten. Diese Daten werden zunächst erfasst, anschließend verarbeitet und schließlich an die Steuerungsebene übergeben.\\\\
In Abbildung \ref{fig:systemaufbau} ist der grundlegende Aufbau des Systems dargestellt. Der Signalfluss beginnt bei der Eingabe, beispielsweise durch einen Scanner, und verläuft über die Verarbeitungsebene, welche durch eine Programmanwendung (z.B. Python) realisiert wird. Innerhalb dieser Ebene erfolgt auch die Synchronisation mit einer Datenbank sowie die Visualisierung über ein HMI. Anschließend werden die verarbeiteten Daten an die SPS übertragen, welche den physikalischen Teil des Prozesses übernimmt.
\begin{figure}[H]
	\centering
	\begin{tikzpicture}[
		node distance=2cm,
		>=latex,
		font=\small,
		every node/.style={align=center}
		]
		
		% Nodes
		\node (scanner) [draw, rectangle, fill=blue!10, minimum width=3.5cm, minimum height=1cm] 
		{Input\\Scanner (RS232/USB)};
		
		\node (python) [draw, rectangle, fill=green!10, below=of scanner, minimum width=4cm, minimum height=1cm] 
		{Programmanwendung\\(Python)\\Datenverarbeitung};
		
		\node (db) [draw, rectangle, fill=yellow!20, right=4cm of python, minimum width=3.5cm, minimum height=1cm] 
		{Datenbank\\(Aufträge / Artikel)};
		
		\node (sps) [draw, rectangle, fill=orange!10, below=of python, minimum width=3.5cm, minimum height=1cm] 
		{SPS\\(z.\,B. Siemens S7-1500)};
		
		\node (actuator) [draw, rectangle, fill=red!10, below=of sps, minimum width=3.5cm, minimum height=1cm] 
		{Aktorik\\(Motor / Lagerbewegung)};
		
		\node (hmi_sps) [draw, rectangle, fill=gray!15, right=4cm of sps, minimum width=3cm, minimum height=1cm] 
		{HMI\\(SPS-seitig)};
		
		% Verbindungen
		\draw[->, thick] (scanner) -- (python) node[midway, right]{RS232};
		
		\draw[->, thick] (python) -- (sps) node[midway, right]{Snap7 / TCP};
		\draw[->, thick] (sps) -- (python);
		
		\draw[->, thick] (sps) -- (actuator) node[midway, right]{Digital I/O};
		
		\draw[->, thick] (sps) -- (hmi_sps) node[midway, above]{PROFINET};
		
		\draw[->, thick] (python) -- (db) node[midway, above]{Abgleich};
		\draw[->, thick] (db) -- (python);
		
	\end{tikzpicture}
	\caption{Systemaufbau der auftragsbasierten Lagersteuerung}
	\label{fig:systemaufbau}
\end{figure}

Wie in Abbildung \ref{fig:systemaufbau} ersichtlich, erfolgt der Ablauf weitgehend seriell. Das bedeutet, dass die einzelnen Prozessschritte nacheinander durchlaufen werden und die Daten schrittweise von der Eingabe bis zur SPS weitergegeben werden. Die SPS stellt dabei die zentrale Schnittstelle zwischen der softwarebasierten Verarbeitungsebene und der physikalischen Anlagenebene dar.\\

Im ersten Schritt erfolgt die Eingabe eines Bauteilcodes, beispielsweise durch das Scannen eines Barcodes. Diese Eingabe wird von der Programmanwendung erfasst und weiterverarbeitet.\\

In der Verarbeitungsebene werden die eingelesenen Daten bereinigt und interpretiert. Auf Basis des Codes erfolgt ein Abgleich mit den in der Datenbank hinterlegten Informationen. Dadurch kann bestimmt werden, welches Bauteil ausgelagert werden soll und welche Lagerposition zugeordnet ist.\\

Nach erfolgreicher Verarbeitung werden die benötigten Steuerinformationen über eine geeignete Kommunikationsschnittstelle an die SPS übertragen. Hierbei kommen je nach System beispielsweise Protokolle wie Snap7 oder ADS zum Einsatz.\\

Die SPS übernimmt anschließend die Ansteuerung der Aktorik. Diese führt die physikalische Bewegung innerhalb des Lagersystems aus, beispielsweise das Verfahren eines Motors oder das Bewegen eines Zylinders zur Auslagerung des gewünschten Bauteils.\\

Die folgenden Abschnitte orientieren sich am in Abbildung\ref{fig:systemaufbau} dargestellten Systemaufbau und beschreiben die einzelnen Komponenten sowie deren Funktion innerhalb des Gesamtsystems.
